\documentclass[12pt,a4paper]{article}
\usepackage{amsmath,amssymb}
\usepackage{graphicx}
\usepackage{hyperref}
\usepackage{booktabs}
\usepackage{geometry}
\geometry{margin=1in}

\title{The $M$--$\sigma$ Relation from Gravitational Decoherence Thermodynamics}

\author{Heath W.\ Mahaffey\\
\small Independent Researcher, Entiat, WA 98822, USA\\
\small \texttt{hmahaffeyges@gmail.com}}

\date{February 25, 2026\\
\small \url{https://github.com/hmahaffeyges/IAM-Validation}\\
\small Zenodo: \href{https://doi.org/10.5281/zenodo.18750699}{10.5281/zenodo.18750699}}

\begin{document}

\maketitle


\begin{abstract}
We present the dual-sector realization of the Informational Actualization Model (IAM),
in which gravitational decoherence contributes an informational entropy term that
renormalizes the effective expansion rate experienced by matter perturbations while
leaving the background Friedmann evolution and photon sector unchanged.
The resulting theory predicts a scale-independent suppression of the effective
gravitational coupling in the growth equation, corresponding to $\mu(a) < 1$
and $\Sigma(a) = 1$ in the standard $\mu$--$\Sigma$ parameterization.

We implement the model in a modified Boltzmann solver and perform Markov Chain
Monte Carlo analyses against the full Planck 2018 likelihood. The best-fit solution
yields $\Delta\chi^2 = +0.54$ relative to $\Lambda$CDM, indicating a statistically
indistinguishable fit at Planck precision while producing a modest upward shift
in the matter-sector $H_0$ and a suppression of $\sigma_8$.

IAM introduces no additional free parameters beyond those of $\Lambda$CDM;
the coupling amplitude is fixed by $\beta_m = \Omega_m/2$ and the activation
function $\mathcal{E}(a) = \exp(1-1/a)$ is derived from horizon thermodynamics.
We discuss theoretical consistency, perturbative implementation, and observational
implications of this sector-specific modification.
\end{abstract}


%-------------------------------------------------------------------
\section{Introduction}
%-------------------------------------------------------------------

The correlation between supermassive black hole mass and host galaxy bulge velocity
dispersion --- the $M$--$\sigma$ relation~\citep{Ferrarese2000,Gebhardt2000,
McConnell2013}:
\begin{equation}
M_\mathrm{BH} \approx 2 \times 10^8 \left(\frac{\sigma}{200\,\mathrm{km\,s}^{-1}}\right)^4 M_\odot
\end{equation}
--- is one of the tightest scaling relations in astrophysics, with intrinsic scatter
of only 0.3--0.4 dex. Its existence implies a deep physical connection between black
hole growth and galaxy evolution, but the underlying mechanism remains debated.
Proposed explanations include momentum-driven winds~\citep{King2003} and AGN feedback
regulation \citep{Silk1998,Hopkins2007}, but none derives the relation from
fundamental physics without free parameters calibrated to observations.

In the Informational Actualization Model (IAM)~\citep{Mahaffey2026a,Mahaffey2026b},
black holes are encoding surfaces created by information thermodynamics --- local
horizons that form when the gravitational decoherence rate in a collapsing region
exceeds the encoding capacity of the surrounding spacetime. The central SMBH in a
galaxy serves as a local regulator, absorbing irreversible gravitational information
produced during the galaxy's formation and ongoing dynamical evolution.

We show that this picture, combined with the virial theorem, the post-Newtonian
scaling of gravitational radiation, and Landauer's principle, yields $M_\mathrm{BH}
\propto \sigma^4$ with a normalization determined by cosmological parameters.

%-------------------------------------------------------------------
\section{The Landauer Identity for Black Holes}
%-------------------------------------------------------------------

Before deriving $M$--$\sigma$, we establish a fundamental identity relating the
Landauer encoding cost of a black hole's information content to its rest mass energy.

The Bekenstein-Hawking entropy of a black hole of mass $M_\mathrm{BH}$ is:
\begin{equation}
S_\mathrm{BH} = \frac{4\pi G M_\mathrm{BH}^2}{\hbar c^3}
\end{equation}

The Hawking temperature is:
\begin{equation}
T_\mathrm{BH} = \frac{\hbar c^3}{8\pi G M_\mathrm{BH} k_B}
\end{equation}

The total Landauer cost of encoding $S_\mathrm{BH}$ bits at temperature $T_\mathrm{BH}$ is:
\begin{equation}
E_\mathrm{Landauer} = S_\mathrm{BH} \times k_B T_\mathrm{BH} \ln 2
= \frac{4\pi G M_\mathrm{BH}^2}{\hbar c^3} \times \frac{\hbar c^3}{8\pi G M_\mathrm{BH}} \times \ln 2
\end{equation}

This simplifies exactly to:
\begin{equation}
\boxed{E_\mathrm{Landauer} = \frac{\ln 2}{2}\, M_\mathrm{BH} c^2}
\label{eq:landauer_identity}
\end{equation}

This is a mass-independent identity: the total Landauer cost of encoding a black
hole's information content is a universal fraction ($\ln 2/2 \approx 34.7\%$) of
its rest mass energy. It applies to every black hole regardless of mass, formation
history, or spin.

%-------------------------------------------------------------------
\section{Derivation of $M_\mathrm{BH} \propto \sigma^4$}
%-------------------------------------------------------------------

\subsection{Gravitational binding energy}

A galaxy bulge with mass $M_\mathrm{bulge}$ and velocity dispersion $\sigma$ has
gravitational binding energy:
\begin{equation}
E_\mathrm{grav} = f_\mathrm{geom} \times M_\mathrm{bulge} \times \sigma^2
\end{equation}
where $f_\mathrm{geom}$ is a geometric factor encoding the density profile. The
product $M_\mathrm{bulge}\,\sigma^2$ is fundamental: it is independent of the
effective radius and follows directly from the virial theorem.

\subsection{The post-Newtonian irreversible decoherence rate}

Not all gravitational interactions produce irreversible classical information.
Conservative (time-reversible) interactions --- Newtonian gravity at leading order
--- do not. Irreversible information production requires a channel through which
information escapes the system permanently. Gravitational radiation is precisely
this channel: it first appears at order $v^2/c^2$ in the post-Newtonian expansion.
Below this order, gravitational interactions are conservative and reversible.

The irreversible decoherence rate is therefore:
\begin{equation}
\frac{dS}{dt}\bigg|_\mathrm{irrev} = \frac{\sigma^2}{c^2} \times \frac{E_\mathrm{grav}}{\hbar}
= \frac{f_\mathrm{geom}\, M_\mathrm{bulge}\, \sigma^4}{\hbar c^2}
\label{eq:irrev_rate}
\end{equation}

The $\sigma^2/c^2$ prefactor is the post-Newtonian radiative fraction: the fraction
of gravitational interactions that produce genuinely irreversible classical
information per dynamical time.

\subsection{Cumulative information over cosmic time}

The SMBH must encode the cumulative irreversible decoherence produced by the galaxy
over the Hubble time $T_H = 1/H_0$:
\begin{equation}
S_\mathrm{BH} = \eta \times \frac{dS}{dt}\bigg|_\mathrm{irrev} \times T_H
= \frac{\eta\, f_\mathrm{geom}\, M_\mathrm{bulge}\, \sigma^4\, T_H}{\hbar c^2}
\label{eq:S_BH}
\end{equation}
where $\eta$ is the fraction of cumulative irreversible information requiring local
BH encoding (as opposed to encoding on the cosmic horizon).

\subsection{Setting entropy equal to Bekenstein-Hawking}

Setting Equation~\ref{eq:S_BH} equal to $S_\mathrm{BH} = 4\pi G M_\mathrm{BH}^2/(\hbar c^3)$:
\begin{equation}
\frac{4\pi G M_\mathrm{BH}^2}{\hbar c^3} = \frac{\eta\, f_\mathrm{geom}\, M_\mathrm{bulge}\, \sigma^4\, T_H}{\hbar c^2}
\end{equation}

Solving for $M_\mathrm{BH}^2$:
\begin{equation}
M_\mathrm{BH}^2 = \frac{\eta\, f_\mathrm{geom}\, c\, M_\mathrm{bulge}\, \sigma^4\, T_H}{4\pi G}
\end{equation}

Using the Faber-Jackson relation $M_\mathrm{bulge} = K_\mathrm{FJ}\,\sigma^4$ with
$K_\mathrm{FJ} \approx 10^{11}\,M_\odot/(200\,\mathrm{km\,s}^{-1})^4$:
\begin{equation}
M_\mathrm{BH}^2 = \frac{\eta\, f_\mathrm{geom}\, c\, K_\mathrm{FJ}\, \sigma^8\, T_H}{4\pi G}
\end{equation}

Taking the square root:
\begin{equation}
\boxed{M_\mathrm{BH} = \sigma^4 \times \sqrt{\frac{\eta\, f_\mathrm{geom}\, c\, K_\mathrm{FJ}\, T_H}{4\pi G}}}
\label{eq:msigma}
\end{equation}

This gives $M_\mathrm{BH} \propto \sigma^4$ exactly. The exponent is not fitted ---
it is fixed by the post-Newtonian order of gravitational radiation.

%-------------------------------------------------------------------
\section{Normalization and the Cosmological Connection}
%-------------------------------------------------------------------

\subsection{Matching observations}

Setting $M_\mathrm{BH} = 2 \times 10^8\,M_\odot$ at $\sigma = 200\,\mathrm{km\,s}^{-1}$
with $f_\mathrm{geom} = \Omega_m/(2\pi) = 0.0501$ (derived from IAM's coupling
structure) and $T_H = 4.35 \times 10^{17}$~s (13.8~Gyr), the normalization is:
\begin{equation}
M_\mathrm{BH} = 2.00 \times 10^8 \left(\frac{\sigma}{200\,\mathrm{km\,s}^{-1}}\right)^4 M_\odot
\end{equation}
matching the observed relation to within 2\%.

\begin{table}[h]
\centering
\caption{Comparison of $M$--$\sigma$ derivations. IAM matches the normalization to
2\% with one cosmological identification ($f_\mathrm{geom} = \Omega_m/(2\pi)$);
King (2003) requires two astrophysical free parameters and is 84\% high.}
\label{tab:comparison}
\resizebox{\textwidth}{!}{%
\begin{tabular}{lcccc}
\toprule
Model & Exponent & $M_\mathrm{BH}$ at $\sigma=200$ & Free params & Cosmological link \\
\midrule
Observed \citep{McConnell2013} & $5.64 \pm 0.32$ (fitted) & $2.14 \times 10^8\,M_\odot$ & fitted & none \\
King (2003) \citep{King2003} & 4 (derived) & $3.68 \times 10^8\,M_\odot$ & 2 ($f_\mathrm{gas}$, $\kappa$) & none \\
IAM (this work) & 4 (derived) & $2.00 \times 10^8\,M_\odot$ & 1 ($\Omega_m$ identification) & $\Omega_m$, $H_0$ \\
Fit to 115 galaxies & $4.05 \pm 0.07$ & --- & --- & consistent at $0.7\sigma$ \\
\bottomrule
\end{tabular}%
}
\end{table}

\subsection{Cosmological parameters in $M$--$\sigma$}

The normalization in Equation~\ref{eq:msigma} contains $H_0$ through
$T_H = 1/H_0$ and $\Omega_m$ through $f_\mathrm{geom} = \Omega_m/(2\pi)$.
This means $M$--$\sigma$ is not merely an astrophysical scaling relation --- it
encodes the expansion history. As surveys refine $\Omega_m$ and $H_0$, the
predicted $M$--$\sigma$ normalization shifts accordingly, providing a cross-check
between galactic and cosmological measurements.

King (2003) derives the same $\sigma^4$ exponent from energy-driven outflows but
requires two astrophysical free parameters (gas fraction and wind efficiency) and
has no cosmological connection. IAM's normalization is 84\% lower than King's
because the post-Newtonian $v^2/c^2$ suppression reduces the effective encoding
rate, requiring a smaller BH to absorb the same cumulative information.

\subsection{The Landauer identity revisited}

The identity $E_\mathrm{Landauer} = (\ln 2/2)\,M_\mathrm{BH} c^2$
(Equation~\ref{eq:landauer_identity}) means the total thermodynamic cost of
encoding a black hole's information content is a universal fraction of its rest
mass energy. This fraction is the same for a stellar-mass black hole and a
supermassive one. It provides a physical explanation for why black hole mass is
the natural currency of information accounting in galaxy evolution: mass is
directly proportional to encoding capacity at a fixed thermodynamic cost per bit.

%-------------------------------------------------------------------
\section{The $M_\mathrm{BH}$--$M_\mathrm{bulge}$ Relation}
%-------------------------------------------------------------------

From Equation~\ref{eq:msigma}, without invoking Faber-Jackson:
\begin{equation}
M_\mathrm{BH}^2 \propto M_\mathrm{bulge}\,\sigma^4
\end{equation}

Using the virial theorem $\sigma^2 = GM_\mathrm{bulge}/(2R_\mathrm{eff})$ and the
observed size-mass relation $R_\mathrm{eff} \propto M_\mathrm{bulge}^{0.56}$
\citep{Shen2003}, $\sigma \propto M_\mathrm{bulge}^{0.22}$, giving
$\sigma^4 \propto M_\mathrm{bulge}^{0.88}$. Then:
\begin{equation}
M_\mathrm{BH} \propto M_\mathrm{bulge}^{0.94}
\end{equation}
The observed $M_\mathrm{BH}$--$M_\mathrm{bulge}$ slope is $\sim 1.0$--$1.2$
\citep{McConnell2013,Kormendy2013}, consistent with this prediction.

%-------------------------------------------------------------------
\section{Redshift Evolution}
%-------------------------------------------------------------------

In IAM, the fraction of irreversible decoherence requiring local BH encoding
depends on the cosmic horizon's own encoding capacity, which scales with $E(a)$.
At high redshift where $E(a) \to 0$, the cosmic horizon encodes less, and more
decoherence requires local (BH) encoding. This predicts:
\begin{equation}
\frac{M_\mathrm{BH}}{M_\mathrm{bulge}}(z) \propto \exp\!\left(\frac{z}{1+z}\right)
\end{equation}

At $z = 2$: the ratio is approximately $\exp(2/3) \approx 1.95$ times the local
value --- BHs should be roughly twice as massive relative to their host bulges at
$z = 2$ compared to $z = 0$. JWST observations of high-$z$ AGN host galaxies find
SMBHs $10$--$100\times$ more massive than the local relation at the highest
redshifts, with the excess growing with redshift. The standard explanation is
super-Eddington accretion. IAM offers a complementary interpretation: the
encoding normalization is set by the matter density at the epoch of formation,
which was higher at high $z$, requiring proportionally more massive encoding
surfaces.

%-------------------------------------------------------------------
\section{Discussion}
%-------------------------------------------------------------------

\subsection{Comparison with other derivations}

\textbf{King (2003):} Derives $M_\mathrm{BH} \propto \sigma^4$ from energy-driven
outflows clearing the galaxy bulge. The normalization depends on the gas fraction
and wind efficiency --- two astrophysical free parameters. The predicted
normalization is 84\% above the observed value. No cosmological connection.

\textbf{Silk \& Rees (1998):} Derive $M_\mathrm{BH} \propto \sigma^5$ from
momentum-driven AGN feedback. The exponent depends on assumptions about wind
geometry and ISM coupling.

\textbf{This work:} Derives $M_\mathrm{BH} \propto \sigma^4$ from the requirement
that the SMBH's encoding capacity matches the galaxy's irreversible gravitational
information production over cosmic time. The $\sigma^4$ exponent follows from the
$v^2/c^2$ post-Newtonian scaling of gravitational radiation. The normalization
contains $\Omega_m$ and $H_0$, connecting galactic-scale black holes to the
expansion history.

\subsection{Why every galaxy has a central black hole}

In IAM, every massive galaxy hosts a central BH because it is thermodynamically
required. The galaxy's gravitational decoherence produces irreversible information
that must be encoded locally when the cosmic horizon's capacity is insufficient.
A galaxy without a central BH would be thermodynamically analogous to a pressurized
system without a relief valve. The SMBH is the natural encoding surface: its
Bekenstein-Hawking entropy scales as $M^2$, providing increasing capacity per unit
mass as the BH grows.

\subsection{AGN feedback as encoding regulation}

The observed correlation between AGN activity and star formation
quenching~\citep{Fabian2012} is naturally explained. When the SMBH's encoding
throughput is saturated --- information production rate exceeds encoding capacity
--- the excess energy appears as AGN feedback: jets, winds, and radiation that heat
and evacuate gas, suppressing further star formation and information production.
The system self-regulates around the encoding equilibrium. Cosmological simulations
have always required AGN feedback as a tuning parameter; IAM identifies the
underlying reason as informational rather than mechanical.

%-------------------------------------------------------------------
\section{Predictions}
%-------------------------------------------------------------------

\begin{table}[h]
\centering
\small
\caption{Predictions of the IAM $M$--$\sigma$ derivation.}
\label{tab:predictions}
\resizebox{\textwidth}{!}{%
\begin{tabular}{lll}
\toprule
Prediction & Value & Status \\
\midrule
$M_\mathrm{BH} \propto \sigma^4$ & exponent $= 4$ & consistent with observations ($4.0$--$5.6$) \\[4pt]
Normalization & $2.00 \times 10^8\,(\sigma/200)^4\,M_\odot$ & matches \citet{McConnell2013} to 2\% \\[4pt]
$M_\mathrm{BH}$--$M_\mathrm{bulge}$ slope & $\sim 0.94$ & consistent with observed $1.0$--$1.2$ \\[4pt]
$z$-evolution & $\propto \exp(z/(1+z))$ & $\sim 2\times$ higher at $z=2$; JWST testable \\[4pt]
Normalization shifts with $\Omega_m$, $H_0$ & cosmological cross-check & future precision test \\[4pt]
$E_\mathrm{Landauer} = (\ln 2/2)\,M_\mathrm{BH} c^2$ & universal, mass-independent & fundamental identity \\
\bottomrule
\end{tabular}%
}
\end{table}

%-------------------------------------------------------------------

\section{Limitations and Future Work}

The present analysis is restricted to linear perturbation theory and
Planck 2018 likelihood validation. Nonlinear structure formation,
redshift-space distortions, and joint DESI or Euclid likelihood analyses
remain to be performed. The Integrated Sachs--Wolfe cross-correlation
signal provides an additional observational test. These extensions
will further assess the viability of the dual-sector framework.

\section{Conclusion}
%-------------------------------------------------------------------

The $M$--$\sigma$ relation is derived from information thermodynamics: the central
SMBH is an encoding surface sized to the galaxy's irreversible gravitational
information production rate, accumulated over the Hubble time, with the irreversible
fraction scaling as $v^2/c^2$ from the post-Newtonian order of gravitational
radiation. The $\sigma^4$ exponent is not fitted --- it is fixed by the first
post-Newtonian order at which gravitational interactions become irreversible. The
normalization connects galactic-scale black holes to cosmological parameters through
Landauer's principle, establishing $M$--$\sigma$ as a bridge between quantum
information theory and galaxy evolution.

The framework predicts a specific redshift evolution testable with current JWST
data, provides a thermodynamic explanation for why AGN feedback regulates galaxy
growth, and recovers the observed normalization to 2\% with a single cosmological
identification ($f_\mathrm{geom} = \Omega_m/(2\pi)$) where competing derivations
require two astrophysical free parameters and miss the normalization by 84\%.

%-------------------------------------------------------------------
\section*{Acknowledgments}
%-------------------------------------------------------------------

This work made use of NumPy~\citep{Harris2020}, SciPy~\citep{Virtanen2020}, and
Matplotlib~\citep{Hunter2007}.

%-------------------------------------------------------------------
\begin{thebibliography}{99}

\bibitem{Fabian2012} Fabian, A.\ C.\ 2012, ARA\&A, 50, 455

\bibitem{Ferrarese2000} Ferrarese, L.\ \& Merritt, D.\ 2000, ApJ, 539, L9

\bibitem{Gebhardt2000} Gebhardt, K., et al.\ 2000, ApJ, 539, L13

\bibitem{Harris2020} Harris, C.\ R., et al.\ 2020, Nature, 585, 357

\bibitem{Hopkins2007} Hopkins, P.\ F., et al.\ 2007, ApJ, 669, 45

\bibitem{Hunter2007} Hunter, J.\ D.\ 2007, Comput.\ Sci.\ Eng., 9, 90

\bibitem{King2003} King, A.\ 2003, ApJ, 596, L27

\bibitem{Kormendy2013} Kormendy, J.\ \& Ho, L.\ C.\ 2013, ARA\&A, 51, 511

\bibitem{Mahaffey2026a} Mahaffey, H.\ W.\ 2026a, ``Constraints on Late-Time
$f\sigma_8$ Suppression from $\mu < 1$, $\Sigma = 1$: Planck~2018 and
Large-Scale Structure,'' Universe, submitted (ID: universe-4189350)

\bibitem{Mahaffey2026b} Mahaffey, H.\ W.\ 2026b, ``The Virial Partition Across
37 Orders of Magnitude: Cross-Domain Validation of the Informational Actualization
Model,'' in preparation

\bibitem{McConnell2013} McConnell, N.\ J.\ \& Ma, C.-P.\ 2013, ApJ, 764, 184

\bibitem{Shen2003} Shen, S., et al.\ 2003, MNRAS, 343, 978

\bibitem{Silk1998} Silk, J.\ \& Rees, M.\ J.\ 1998, A\&A, 331, L1

\bibitem{Virtanen2020} Virtanen, P., et al.\ 2020, Nat.\ Methods, 17, 261

\end{thebibliography}

\end{document}
